% robotaction.tex: Robot retargeting sections for EgoDelta (CoRL 2026)
% Main paper subsection plus appendix detail

%% MAIN PAPER

\subsection{Retargeting 3D Change Targets to Robot Actions}
\label{subsec:retargeting}

Given an actionable 3D change target $g_{\Delta}$, such as a fruit
contact region or a stem/leaf cut point, the retargeting module first
associates it with a metric point $\mathbf{p}_s$ in a calibrated
sensor frame $s$.  For Aria-only targets, $s=\mathrm{aria}$; for live
RGB-D execution, $s=\mathrm{rs}$ after EgoDelta instance association
and RealSense depth lookup.

In this benchmark, harvesting GT marks the tomato apex on the removed
fruit-and-calyx side of the interface with the retained green truss,
whereas deleafing GT marks the magnetic target on the retained stem-rod
side of the petiole interface.
We therefore represent each attachment by the interface pair
$(g_{\mathrm{ret}},g_{\mathrm{rem}})$ and let the downstream tool/action
convention select the mechanically appropriate endpoint.

With $\bar{\mathbf{p}}_s=[\mathbf{p}_s^{\top},1]^{\top}$,
%
\begin{equation}
  \mathbf{p}_{\mathrm{robot}}
  = \Pi\!\left(
      \mathbf{T}_{\mathrm{robot}}^{s}\,
      \bar{\mathbf{p}}_s
    \right),
  \label{eq:target_transform}
\end{equation}
%
where $\Pi$ dehomogenizes the transformed point and
$\mathbf{T}_{\mathrm{robot}}^{\mathrm{aria}} =
\mathbf{T}_{\mathrm{robot}}^{\mathrm{rs}}
\mathbf{T}_{\mathrm{rs}}^{\mathrm{aria}}$ when the source frame is
Aria.

When a live RGB-D view is available, we use the EgoDelta mask for
instance identity and RealSense depth for metric localization.  This
keeps the semantic target tied to the egocentric video while
reducing sensitivity to depth error in the Aria reconstruction.

The robot does not replay the human trajectory.  Instead, the
robot-frame target instantiates a task primitive
$\pi(\mathbf{p}_{\mathrm{robot}}, \mathcal{G})$, where
$\mathcal{G}$ contains the current RealSense point cloud, calibrated
tool geometry, and Aria MPS hand-motion priors.  MPS provides
phase-level action structure: approach, contact, and post-contact
directions.  Fixed task templates provide scalar parameters such as
standoff, insertion, lift, and return height.

For harvesting, the primitive turns a fruit target into a gripper
contact frame and target-derived insertion plan.  For deleafing, the
primitive turns a stem/leaf target into a calibrated knife-edge frame
and a near-vertical cut plan with point-cloud-scored entry.  Because
each primitive is instantiated from the current 3D target and local
geometry, the same action template can be applied zero-shot to new
targets of the same action class.  A bounded IK solver converts
candidate end-effector waypoints to arm joint targets under
joint-limit, self-clearance, point-cloud path-clearance, and
task-contact gates.
Appendix~\ref{app:robot_retargeting} gives the full construction.


%% APPENDIX

\subsection{Robot Retargeting}
\label{app:robot_retargeting}

\subsubsection{Target Localization and Task Frames}
\label{app:task_frame}

The retargeting module consumes a robot-frame target
$\mathbf{p}_{\mathrm{robot}}$ and an action type.  The target may be
obtained by projecting an EgoDelta mask into a live RealSense depth
image or supplied as a metric 3D anchor for hardware validation.
Both interfaces produce the same
quantity: a metric 3D anchor in the robot planning frame.

In this benchmark, harvesting GT marks the tomato apex on the removed
fruit-and-calyx side of the interface with the retained green truss,
whereas deleafing GT marks the magnetic target on the retained stem-rod
side of the petiole interface.
We therefore represent each attachment by the interface pair
$(g_{\mathrm{ret}},g_{\mathrm{rem}})$ and let the downstream tool/action
convention select the mechanically appropriate endpoint.

For harvesting, the anchor is expanded into a fruit task frame
containing:
%
\begin{itemize}
  \item a contact point $\mathbf{p}_c$, the midpoint of the planned jaw
    contact locations on the fruit surface;
  \item a fruit center $\mathbf{p}_f$, estimated from the 3D mask
    volume;
  \item an orthonormal triad
    $(\hat{\mathbf{a}},\,\hat{\mathbf{j}},\,\hat{\mathbf{n}})$:
    approach/backoff axis, jaw axis, and palm normal;
  \item scalar fruit width $w$ and depth $d$.
\end{itemize}
%
The jaw axis is initialized from mask contact geometry and may be
refined by the MPS pinch-axis prior when the hand-tracking confidence
is sufficient.

For deleafing, the anchor is the desired cut point
$\mathbf{p}_{\mathrm{cut}}$.  The task frame contains an approach direction
$\hat{\mathbf{a}}$, a cut direction $\hat{\mathbf{s}}$, an entry
side direction $\hat{\mathbf{c}}$, and the calibrated knife-frame axes.
The approach and raw cut directions come from MPS phase motion when
available.  The entry direction is inferred from the local
RealSense point cloud by moving away from nearby occupied points around
the cut target.


\subsubsection{Harvest Primitive}
\label{app:harvest_primitive}

For a side grasp on a roughly spherical fruit, the gripper must advance
past the surface contact point along $-\hat{\mathbf{a}}$ so that the
jaws reach the widest cross-section.  We define the grasp anchor
%
\begin{equation}
  \mathbf{p}_g
  = \mathbf{p}_c
    - \delta_{\mathrm{ins}}\,\hat{\mathbf{a}}
    + \delta_{\mathrm{palm}}\,\hat{\mathbf{n}},
  \label{eq:grasp_anchor}
\end{equation}
%
where $\delta_{\mathrm{ins}}$ is the insertion depth and
$\delta_{\mathrm{palm}}$ is a bounded palm-axis correction.  The
nominal insertion is $\delta_{\mathrm{ins}}^{*}=d/2$.  The planner
sweeps insertion factors around this nominal value and records the
solved contact position, not only the requested target, before
accepting a grasp.

Harvesting uses the sequence
\emph{approach} $\rightarrow$ \emph{pre-grasp} $\rightarrow$
\emph{grasp} $\rightarrow$ \emph{harvest} $\rightarrow$
\emph{retract} $\rightarrow$ \emph{basket}.  The
approach-through-grasp sub-sequence is solved first and then frozen;
post-grasp waypoints cannot change the validated grasp contact.  At
pre-grasp the gripper opens, at grasp it closes, and at basket it
releases.  Post-grasp waypoints use the MPS pull direction when
confident, with a small downward detach component for stem separation.


\subsubsection{Deleafing Primitive}
\label{app:deleaf_primitive}

Deleafing is modeled as a short blade stroke through a 3D cut point.
The current implementation uses a fixed method configuration rather
than per-target tuning.  Given a cut point
$\mathbf{p}_{\mathrm{cut}}$, the planner offsets the calibrated blade
by a configured target-extension distance along
$\hat{\mathbf{u}}_{\mathrm{engage}}$, the horizontal target-extension
direction for the blade.  The cut waypoint then moves an additional
distance $d_{\mathrm{cut}}$ along the stroke axis.  The stroke
direction is constrained to lie within $\theta_{\max}$ of vertical:
if the MPS cut axis already satisfies this bound after being flipped
toward world down, it is preserved; otherwise, its horizontal
component sets the clamped tilt direction.  Degenerate horizontal
motion falls back to world down.  This preserves the human-indicated
stroke side while enforcing a near-vertical cut that reduces lateral
wrist sweep near the plant support rods.

The deleafing waypoint sequence is
%
\begin{enumerate}
  \item \emph{Transit lift}: lift near the home side before entering
    the plant workspace.
  \item \emph{Entry transit}: approach the local target from a
    point-cloud-selected side direction.
  \item \emph{Approach}: align the knife above the cut line.
  \item \emph{Pre-cut}: stop at a fixed standoff along
    $-\hat{\mathbf{s}}$.
  \item \emph{Cut}: move a few millimeters through the target along
    $\hat{\mathbf{s}}$.
  \item \emph{Post-cut lift}: lift back along $-\hat{\mathbf{s}}$.
  \item \emph{Return transit}: move high toward home before the
    standard home restore.
\end{enumerate}
%
Candidate generation varies the entry-side angle and contact
advance distance using the current RealSense point cloud.  A candidate
is accepted only if all requested waypoints solve, interpolated
self-clearance passes, and local point-cloud clearance passes up to
the requested execution waypoint.

All deleafing waypoints are recomputed from the current cut point,
MPS-derived action axes, method-configured tool orientation,
calibrated knife geometry, and the live RealSense point cloud.  The
method therefore changes target position and local entry side without
per-point parameter tuning.


\subsubsection{Inverse Kinematics}
\label{app:ik}

Each waypoint is solved as a bounded nonlinear least-squares problem
over the active arm joints $\mathbf{q}\in\mathbb{R}^7$.  The base,
torso, head, and inactive arm remain fixed at the home configuration.

\paragraph{Full-pose objective.}
For waypoints requiring a full TCP orientation, the residual is
%
\begin{equation}
  \mathbf{r}(\mathbf{q})
  =
  \begin{bmatrix}
    w_p(\mathbf{t}(\mathbf{q})-\mathbf{t}^*) \\
    w_o\,\mathrm{Log}(\mathbf{R}^{*\top}\mathbf{R}(\mathbf{q})) \\
    w_h(\mathbf{q}-\mathbf{q}_{\mathrm{home}}) \\
    w_s(\mathbf{q}-\mathbf{q}_{\mathrm{ref}}) \\
    w_c\,\mathbf{c}(\mathbf{q})
  \end{bmatrix},
  \label{eq:ik_residual}
\end{equation}
%
where $\mathbf{t}$ and $\mathbf{R}$ are the calibrated TCP position
and rotation, $\mathrm{Log}$ is the $\mathrm{SO}(3)$ logarithm, and
$\mathbf{c}$ contains one-sided self-clearance penalties for the
nearest non-adjacent link pairs.  The vector
$\mathbf{q}_{\mathrm{home}}$ is the stored home posture, while
$\mathbf{q}_{\mathrm{ref}}$ is the previous waypoint solution.

\paragraph{Task-axis objective.}
For contacts where one wrist degree of freedom is mechanically
irrelevant, we replace the orientation residual by axis alignment:
%
\begin{equation}
  \mathbf{r}_{\mathrm{ori}}(\mathbf{q})
  =
  \begin{bmatrix}
    w_1\,\hat{\mathbf{x}}_1^* \times \hat{\mathbf{x}}_1(\mathbf{q}) \\
    w_2\,\hat{\mathbf{x}}_2^* \times \hat{\mathbf{x}}_2(\mathbf{q})
  \end{bmatrix}.
  \label{eq:task_axis_residual}
\end{equation}
%
For harvesting, the axes are the gripper jaw and approach directions.
For deleafing, they are the calibrated knife edge and backoff
directions.  The local tool axes may be canonical TCP axes or
calibrated vectors, which allows the same solver to handle the
parallel-jaw gripper and the right-arm knife.

\paragraph{Solver and seeding.}
We minimize $\|\mathbf{r}(\mathbf{q})\|^2$ with the Trust Region
Reflective algorithm~\cite{branch1999trf} under joint-position bounds.
Each waypoint is seeded from the previous solved waypoint, which
encourages continuity while still solving the full IK problem for the
current target.

\paragraph{Tool calibration.}
The TCP is
$\mathbf{T}_{\mathrm{tcp}}=\mathbf{T}_{\mathrm{eef}}(\mathbf{q})
\mathbf{T}_{\mathrm{eef}}^{\mathrm{tcp}}$.  Harvesting further uses a
gripper contact-center offset.  Deleafing uses a right-knife pivot
calibration for the tool tip and an edge-line calibration for the
knife cutting direction.  Thus the IK target is expressed in the
physical contact frame of the tool, not merely at the wrist.


\subsubsection{Candidate Scoring and Safety Gates}
\label{app:gates}

For each primitive, the planner generates a finite candidate set and
evaluates it with the same gates used by the executor:
%
\begin{itemize}
  \item \textbf{IK convergence}: waypoint position error below the
    task tolerance and orientation or task-axis error below the action
    tolerance.
  \item \textbf{Self-clearance}: endpoint and interpolated joint paths
    must clear non-adjacent robot links by the configured margin.
  \item \textbf{Point-cloud clearance}: harvest candidates check
    modeled non-target fruit and gripper volume; deleaf candidates
    check the TCP path against the live RealSense point cloud near the
    cut point.
  \item \textbf{Contact validity}: harvest candidates must satisfy
    open-gripper alignment and solved insertion depth; deleaf
    candidates must preserve the fixed cut standoff, target extension,
    and downward stroke.
  \item \textbf{Joint continuity}: consecutive waypoint joint changes
    are bounded to avoid wrap-around solutions.
\end{itemize}
%
Candidates are ranked lexicographically by gate success, point-cloud
margin, self-clearance margin, and endpoint error.  For deleafing,
ties are further broken by agreement with the configured target
extension and smaller contact advance.
All candidates and gate evaluations are recorded per trial for
analysis, including selected waypoints, IK errors, clearance reports,
MPS axes, tool calibration, and the execution decision.


\subsubsection{MPS Priors and Task Templates}
\label{app:hand_prior}

MPS hand tracking is used as an action prior, not as a trajectory to
replay.  From the egocentric video we extract phase windows and compute
median hand-motion axes in the robot frame.  Harvesting uses
approach, jaw/pinch, and pull axes.  Deleafing uses approach, cut, and
retract axes from the wrist or pinch trajectory.  Each axis carries a
confidence derived from sample count and motion magnitude; low
confidence axes fall back to mask geometry, world-down stroke, or
home-to-target approach directions.

Task templates define target-independent action quantities, including
standoff, target extension, stroke length, return lift, and
tool-orientation constraints.  At runtime, the planner instantiates
the MPS-guided template at the current 3D target and solves a fresh
waypoint sequence in the robot frame.


\subsubsection{Execution and Safety}
\label{app:execution}

Hardware execution is separated from planning by a hard gate: commands
are issued only after the selected candidate passes all requested
waypoint checks and operator confirmation is provided.  Before motion,
the executor verifies the live robot state against the stored home
configuration.  Waypoints are sent as time-parameterized joint targets;
gripper commands are issued only at validated grasp/release waypoints.

In the default deleafing configuration, execution stops at the
return-transit waypoint.  The robot is then restored with an arm-only
home command, so the knife does not drag downward through the plant
while returning from the post-cut state.
